\section{Conclusions}
\label{sec:Conclusion}
In this work an active D band balun with low amplitude and phase imbalance in SG13G3 \SI{130}{\nano\meter}-Sige BiCMOS techonolgy was presented.\\
  
The designed balun with simultaneously terminated outputs presents a mismatch of \SI{\pm 0.2}{\decibel} and \SI[separate-uncertainty=true]{-1 \pm 1}{\degree} in the \SIrange{100}{200}{\giga\hertz} range.
The \SI{3}{\decibel} gain bandwidth is achieved in the \SIrange{112}{175}{\giga\hertz} range, with a maximum gain of \SI{12.63}{\decibel} at \SI{112}{\giga\hertz} and a minimum gain of \SI{9.63}{\decibel} at \SI{140}{\giga\hertz} and \SI{175}{\giga\hertz}.\\

Results obtained by simulating one electromagnetic termination at a time are presented for measurement purposes. Over the \SIrange{100}{200}{\giga\hertz} bandwidth, the balun shows an amplitude imbalance of \SI[separate-uncertainty=true]{-0.1 \pm 1}{\decibel} and a phase imbalance of \SI[separate-uncertainty=true]{1 \pm 5}{\degree}. This discrepancy is caused by unbalanced single-ended terminations, which generate mixed-mode conversion.\\

The simulations and the comparisons with EM cores provided in section (\ref{Analysis_combined_core}) indicate that good balun performance, in terms of amplitude and phase imbalance, can be achieved using only the first two differential stages. This suggests that the amplifier stage may be replaced if different gain performance is required.\\

The results are based solely on simulations, since the circuit was submitted for tape-out only at a later stage.
\begin{comment}
    An alternative to the differential topologies is also presented: A first schematic and a carefull layout realization of the CBCE balun is capable of achieving amplitude balance of \SI{0.4}{\decibel}@\SI{180}{\giga\hertz} by itself. altought the phase imbalance needs more precautions. Further studies are possible.\\
\end{comment}


